\begin{center}
\textbf{\large Tóm tắt}
\end{center}
\addcontentsline{toc}{chapter}{Tóm tắt}
\begin{small}

Một hàm Python có thể vượt hết test cơ sở nhưng trả sai khi gặp một đầu vào khác. Khi điểm trên benchmark sinh mã tiến gần mức tối đa, những trường hợp bị bỏ sót như vậy làm giảm khả năng phân biệt các lời giải. ConceptTesting hướng dẫn mô hình ngôn ngữ lớn (LLM) chọn đầu vào kiểm thử. Phân hoạch tương đương kết hợp phân tích giá trị biên (EP/BVA) chọn đại diện của lớp và giá trị sát biên, còn bảng quyết định (DT) chọn tổ hợp điều kiện. Hai nhánh dùng phản hồi từ lượt đầu để bổ sung mục tiêu còn thiếu. ControlFlow dùng dòng, cạnh chưa phủ và nghĩa vụ theo đặc tả để chọn tình huống cần thử trên HumanEval. Thực nghiệm dùng hai bộ dữ liệu chính là HumanEval và MBPP, với đáp án lấy từ các cài đặt tham chiếu của EvalPlus.

Trên HumanEval, EP/BVA và DT hai lượt với Qwen phát hiện sai khác trên 154 trong 1.454 chương trình được kiểm thử (PUT) đã vượt test cơ sở. ControlFlow theo nghĩa vụ phát hiện thêm 20 PUT ngoài hai nhánh, trong đó 18 liên quan đến giới hạn tài nguyên. Hợp ba nhánh phát hiện 174 PUT, so với 87 PUT khi Qwen sinh test trực tiếp. Pass@1 trên 162 bài đủ mười cài đặt giảm từ 89,75\% xuống 79,01\%. Khi loại mười phát hiện phụ thuộc cách hiểu Boolean ở một bài, điểm còn 79,63\%. Hợp ba nhánh nâng số PUT đạt C2 100\% từ 697 lên 866 trên 1.151 PUT có nhánh. Trên tập 164 bài MBPP+, Qwen sinh test trực tiếp phát hiện sai khác lặp lại ở 274 trong 961 PUT đã vượt test cơ sở và nâng số PUT đạt C2 100\% từ 174 lên 429 trên 634 PUT có nhánh.

\vspace*{.6cm}
\textbf{Từ khóa}\quad sinh ca kiểm thử, khái niệm kiểm thử, phát hiện lỗi, benchmark sinh mã, mô hình ngôn ngữ lớn.

\end{small}
